% 题干转写：完成
% 解析转写：完成
% 图形核对：完成
\section{第38届复赛第4题：两枚小永磁针的相互作用与平衡}

\subsection{题干}
一枚小永磁针可视为半径很小的电流环，其磁矩大小为 $\mu=IS$，其中 $I$ 为环电流强度、$S=\pi R^2$，磁矩方向与电流环平面垂直并满足右手螺旋关系。两枚小磁针 $A$、$B$ 的磁矩大小均为 $\mu$、质量均为 $m$。取竖直向下为 $z$ 轴正向，$A$ 固定在原点 $O$，磁矩沿 $x$ 轴正向；$B$ 位于 $A$ 正下方。已知真空磁导率 $\mu_0$ 和重力加速度大小 $g$。

设 $B$ 的质心坐标为 $(0,0,z)$，且 $z\gg R$；其磁矩与 $x$ 轴夹角为 $\theta$。求：（1）$A$、$B$ 间的相互作用势能；（2）在 $z$ 不变时的稳定平衡角 $\theta_{\text{稳}}$；（3）当 $\theta=\theta_{\text{稳}}$ 而 $z$ 可变化时的受力平衡位置，并判断其稳定性。
\begin{figure}[htbp]
  \centering
  \begin{tikzpicture}[x=.9cm,y=.9cm,>=Stealth,font=\small]
  \coordinate (O) at (0,0);
  \draw[thick] (O) ellipse (1.65 and .54);
  \draw[->,very thick] (-1.15,.25) arc (150:35:1.65 and .54)
    node[pos=.46,above] {$I$};
  \fill (O) circle (1.5pt);
  \draw[densely dashed] (O)--(-.82,.33) node[midway,above] {$R$};
  \draw[->,very thick] (O)--(.7,1.55) node[right] {$\boldsymbol{\mu}$};
  \node[below] at (O) {电流环};

  \begin{scope}[xshift=5.2cm]
    \coordinate (A) at (0,1.3);
    \coordinate (B) at (0,-1.45);
    \fill (A) circle (1.7pt) node[left] {$A,\ O$};
    \fill (B) circle (1.7pt) node[left] {$B$};
    \draw[->,very thick] (A)--(1.35,1.3) node[above] {$\boldsymbol{\mu}_A\parallel +x$};
    \draw[->,very thick] (B)--(1.1,-.68) node[right] {$\boldsymbol{\mu}_B$};
    \draw[densely dashed] (A)--(B) node[midway,left] {$z$};
    \draw[->] (-.85,1.3)--(-.85,-1.45) node[midway,left] {$+z$};
    \draw (B)+(.62,0) arc (0:35:.62) node[midway,right] {$\theta$};
    \node[align=center,below] at (.35,-2.05) {磁偶极子相对取向\\（示意）};
  \end{scope}
\end{tikzpicture}

  \caption{第38届复赛第4题电流环磁矩与两磁偶极子相对位置的TikZ重绘。左图对应原题图4a；右图用于表示题设中的相对取向。}
  \label{fig:cpho38q04}
\end{figure}

\subsection{参考解答整理}
\paragraph{（1）远场磁场与相互作用势能。}
把 $A$ 看作磁偶极子。由毕奥--萨伐尔定律在 $z\gg R$ 下展开并对环积分，$B$ 所在位置的磁场为
\[
\mathbf B_A=-\frac{\mu_0\mu}{4\pi z^3}\,\mathbf e_x.
\]
取 $\boldsymbol\mu_B$ 与磁场垂直时的势能为零，则
\[
\boxed{E_p=-\boldsymbol\mu_B\cdot\mathbf B_A
=\frac{\mu_0\mu^2\cos\theta}{4\pi z^3}.}
\]

\paragraph{（2）磁矩指向的稳定平衡。}
在 $z$ 固定时，重力势能不影响对 $\theta$ 的判定，总势能可写为
\[
V(\theta)=\frac{\mu_0\mu^2\cos\theta}{4\pi z^3}-mgz.
\]
外力矩大小为
\[
M=-\frac{\mathrm dV}{\mathrm d\theta}
=\frac{\mu_0\mu^2\sin\theta}{4\pi z^3}.
\]
因此候选平衡角为 $\theta=0,\pi$。又有
\[
\frac{\mathrm d^2V}{\mathrm d\theta^2}
=-\frac{\mu_0\mu^2\cos\theta}{4\pi z^3}.
\]
在 $\theta=0$ 处二阶导数为负，在 $\theta=\pi$ 处为正，故
\[
\boxed{\theta_{\text{稳}}=\pi.}
\]

\paragraph{（3）竖直方向的受力平衡。}
令 $\theta=\pi$，以竖直向下为 $z$ 正向。由 $F=-\mathrm dV/\mathrm dz$ 得
\[
F=-\frac{3\mu_0\mu^2}{4\pi z^4}+mg.
\]
令 $F=0$，受力平衡位置为
\[
\boxed{z_0=\left(\frac{3\mu_0\mu^2}{4\pi mg}\right)^{1/4}.}
\]
在该处
\[
\frac{\mathrm d^2V}{\mathrm dz^2}
=-\frac{3\mu_0\mu^2}{\pi z_0^5}<0.
\]
故该点为势能极大值，对竖直扰动不稳定：$B$ 下移时重力占优而继续下落，上移时磁吸引占优而趋向 $A$。

\subsection{转写依据}
题干依据 \texttt{CPHO38-Q04\_transcription\_v1.md}；公式依据 \texttt{CPHO38-Q04\_formula\_check\_v1.csv} 与答案第10至12页；左侧电流环依据原题图4a，右侧相对位置图仅用于呈现题设几何关系。本题第（1）、（2）、（3）问评分分别为14分、13分、13分。
